\documentclass[parskip=half-, numbers=enddot]{scrartcl}
\usepackage[T1]{fontenc}
\usepackage[utf8]{inputenc}
\usepackage{mathtools}
\usepackage{cmbright}
\usepackage{microtype}
\usepackage{tabularray}
\UseTblrLibrary{booktabs}

\begin{document}
\begin{center}
  \bfseries
  \begin{Huge}
    Intentions pédagogiques
  \end{Huge}
  \smallbreak
  \begin{Large}
    Mathématiques - 1ères
  \end{Large}
  \smallbreak
  \rule[10pt]{.3\linewidth}{1 pt}
\end{center}
\section{Compétences du cours de mathématiques}

\subsection{Expliciter les savoirs et les procédures}

Il s'agit d'évoquer les connaissances qui s'y rapportent et montrer
qu'on en saisit le sens, la portée. Cela se traduit par des
justifications, des explications, des erreurs à retrouver \dots

\subsection{Appliquer des concepts et des procédures}

Au moment où l'élève apprend une procédure, il opère des raisonnements
et construit des enchaînements qui ne sont pas d'emblée des
automatismes. À son niveau, ces techniques sont parfois complexes. Il
s'agit pour lui d'acquérir des réflexes réfléchis. Il s'agira donc de
réaliser des exercices qui ressemblent fortement aux exercices déjà
rencontrés. La réalisation des exercices nécessite : une connaissance de
diverses propriétés, de règles de calcul, des procédés de constructions
des différentes figures planes, de la représentation d'objets dans
l'espace, une habileté calculatoire, la maîtrise de la
calculatrice \dots

\subsection{Résoudre des problèmes}

La résolution de problèmes place l'élève, soit dans un contexte qui
n'est pas encore mathématisé, soit dans une situation nouvelle. Il devra
utiliser les techniques vues en classe pour les résoudre. Cela demande
d'établir des connexions entre les différents points de la matière.
Enfin, il doit pouvoir commenter, justifier les étapes de son
raisonnement et interpréter les résultats obtenus.

\section{Évaluations}

Tout au long de l'année, l'élève sera amené à évaluer sa progression et
sa compréhension des matières au travers des interrogations. Elles se
distinguent en deux catégories.

\section{Évaluations régulières ( Travail journalier - TJ)}

Cela représente toutes les interrogations effectuées en cours de
chapitre. Elles peuvent concerner un point de matière particulier, des
exercices de drill, une vérification de l'acquisition de la théorie du chapitre, plusieurs points de matière \dots Certaines peuvent être
programmées et d'autres non.

\section{Bilans de chapitres (B)}

Il s'agit des interrogations qui reprennent l'ensemble des matières d'un chapitre ou de plusieurs chapitres.

\section{Pondération}

L'année est évaluée sur un total de 100 points. \\
La pondération du cours de mathématiques se fait selon la règle suivante: 
\begin{itemize}
  \item Toutes les périodes portent sur 7 points, à raison de 2 points pour le TJ et 5 points pour les bilans.
  \item L'examen de Noël porte sur 30 points.
  \item L'examen de juin porte sur 35 points.
\end{itemize}
 
En résumé :
\begin{center}
  \begin{tblr}{row{1}={font=\bfseries}}
    \toprule
    P1 & P2 & Noël & P3 & P4 & P5 & Juin & Total \\
    \midrule
    7 & 7 & 30   & 7 & 7 & 7 & 35 & 100 \\
    \bottomrule
  \end{tblr}
\end{center}

\section{Dossier d'apprentissage}

\subsection{Importance du dossier d'apprentissage}

Les dossiers d'apprentissage sont des outils précieux qui font partie
des documents officiels à conserver pour l'inspection et l'homologation.
Pour chacune des matières, le dossier (farde à devis) contient toutes
les évaluations de l'élève, lui permettant de se situer dans ses
apprentissages et permettant à ses parents de suivre le travail de leur
enfant. C'est en analysant et en corrigeant les erreurs faites lors des
évaluations, que l'élève pourra progresser. Il est donc important que
chacun de ces dossiers soit bien tenu.

\subsection{Contenu et tenue du dossier d'apprentissage}

Les évaluations et les devoirs cotés seront corrigés et mis dans le
dossier d'apprentissage. Ils sont classés par ordre chronologique (du
plus récent au plus ancien) et la feuille récapitulative, appelée
répertoire, sera placée comme page de garde. Les devoirs et les
évaluations (régulières ou bilans) seront faits sur feuille à en-tête,
agrafées si nécessaire (sauf si les élèves répondent sur le
questionnaire). En début d'année, chaque élève reçoit un bloc de
feuilles à en-tête de l'Institut. Des blocs supplémentaires sont en
vente à l'accueil.

\section{Critère de réussite}

Le critère de réussite du cours de mathématiques est d'obtenir une note dépassant ou égalant les 50\% au global de l'année. 
Attention : Un travail de remise à niveau pourra toutefois être donné fin juin aux élèves qui ont réussi l'année mais qui ont obtenu moins de 50\% à l'examen de fin d'année.

\vfil
Signature de l'élève
\hfill
Signature d'un responsable légal

\end{document}
